% Appendix B
\section{Exercises, with Solutions}\label{app:exercises}

Exercises are ordered by chapter. Each solution is a miniature version of an argument in the text — do them with the book closed, then read.

\subsection*{Chapter \ref{sec:ls}}

\begin{exercise}\label{ex:sigma}
Show that $\hat{\sigma}^2 = \norm{\bm{e}}^2/(n-p-1)$ is unbiased for $\sigma^2$ under (A1)--(A3), i.e.\ $\E[\norm{\bm{e}}^2] = (n-p-1)\sigma^2$.
\end{exercise}
\begin{proof}
Write $\bm{e} = (\bm{I} - \bm{H})\y = (\bm{I}-\bm{H})(\X\bbeta + \bm{\varepsilon}) = (\bm{I}-\bm{H})\bm{\varepsilon}$, since $(\bm{I}-\bm{H})\X = \zero$. With $\bm{M} := \bm{I} - \bm{H}$ (symmetric, idempotent):
\[
  \E[\norm{\bm{e}}^2] = \E[\bm{\varepsilon}^\top \bm{M} \bm{\varepsilon}] = \E[\trace(\bm{\varepsilon}^\top \bm{M} \bm{\varepsilon})] = \E[\trace(\bm{M}\bm{\varepsilon}\bm{\varepsilon}^\top)] = \trace\big(\bm{M}\, \E[\bm{\varepsilon}\bm{\varepsilon}^\top]\big) = \sigma^2 \trace(\bm{M}) = \sigma^2 (n - p - 1),
\]
using $\E[\bm{\varepsilon}\bm{\varepsilon}^\top] = \Var(\bm{\varepsilon}) = \sigma^2 \bm{I}$ and Lemma~\ref{lem:trace}. Hence $\E[\hat{\sigma}^2] = \sigma^2$.
\end{proof}

\begin{exercise}\label{ex:means}
In simple regression (one feature), verify directly that the least-squares line passes through $(\bar{x}, \bar{y})$ and that \eqref{eq:slope} satisfies the second normal equation.
\end{exercise}
\begin{proof}
The intercept equation gives $\hat\beta_0 = \bar{y} - \hat\beta_1 \bar{x}$, so the fitted value at $\bar{x}$ is $\hat\beta_0 + \hat\beta_1 \bar{x} = \bar{y}$. For the slope equation, substitute $\hat\beta_0$:
\[
  \sum_i x_i\big( (y_i - \bar{y}) - \hat\beta_1 (x_i - \bar{x}) \big) = \sum_i (x_i - \bar{x})(y_i - \bar{y}) - \hat\beta_1 \sum_i (x_i - \bar{x})^2 = 0,
\]
using $\sum_i (x_i - \bar{x}) = 0$ to replace $x_i$ by $x_i - \bar{x}$, and the definition \eqref{eq:slope} of $\hat\beta_1$.
\end{proof}

\subsection*{Chapter \ref{sec:why-squared}}

\begin{exercise}
Show that minimizing $\sum_i |y_i - \theta|$ over $\theta$ yields the sample median. What goes wrong with the derivative-based approach, and how does subgradient calculus fix it?
\end{exercise}
\begin{proof}
$|e|$ is not differentiable at $e = 0$; its subgradient is $\mathrm{sign}(e)$ for $e \neq 0$ and $[-1, 1]$ at $e = 0$. Zero must lie in the subdifferential of the objective: $0 \in \sum_i \partial|y_i - \theta|$. For $\theta$ not equal to any observation, the sum equals $\#\{y_i > \theta\} - \#\{y_i < \theta\}$, which is zero exactly when $\theta$ has as many observations above as below — the median. Between observations the subgradient is constant, so the objective is piecewise linear and any median attains the minimum.
\end{proof}

\subsection*{Chapter \ref{sec:geometry}}

\begin{exercise}
Prove that for any $\bm{y} \in \R^n$, $\norm{\bm{H}\bm{y}} \leq \norm{\bm{y}}$, with equality iff $\bm{y} \in \mathcal{C}(\X)$. Conclude $0 \leq H_{ii} \leq 1$.
\end{exercise}
\begin{proof}
$\norm{\bm{y}}^2 = \norm{\bm{H}\bm{y}}^2 + \norm{(\bm{I}-\bm{H})\bm{y}}^2 \geq \norm{\bm{H}\bm{y}}^2$ by orthogonality; equality iff $(\bm{I}-\bm{H})\bm{y} = \zero$, i.e.\ $\bm{y} = \bm{H}\bm{y} \in \mathcal{C}(\bm{H}) = \mathcal{C}(\X)$. For the diagonal: $H_{ii} = \bm{e}_i^\top \bm{H}\bm{e}_i = \bm{e}_i^\top \bm{H}^2 \bm{e}_i = \norm{\bm{H}\bm{e}_i}^2 \geq 0$ and $\norm{\bm{H}\bm{e}_i} \leq \norm{\bm{e}_i} = 1$ by the first part (with $\bm{H}\bm{e}_i \in \mathcal{C}(\X)$, so in fact $H_{ii} = \norm{\bm{H}\bm{e}_i}^2 \in [0,1]$).
\end{proof}

\subsection*{Chapter \ref{sec:gauss-markov}}

\begin{exercise}\label{ex:gls}
Let $\Var(\bm{\varepsilon}\mid\X) = \bm{\Sigma}$ with $\bm{\Sigma} \succ \zero$ known. Show that $\bhat_{\mathrm{GLS}} = (\X^\top\bm{\Sigma}^{-1}\X)^{-1}\X^\top\bm{\Sigma}^{-1}\y$ is the BLUE of $\bbeta$ by applying Gauss--Markov in whitened coordinates.
\end{exercise}
\begin{proof}
The symmetric positive definite $\bm{\Sigma}$ has a square root $\bm{\Sigma}^{1/2} \succ \zero$ (eigendecomposition). Multiply the model by $\bm{\Sigma}^{-1/2}$: $\bm{\Sigma}^{-1/2}\y = \bm{\Sigma}^{-1/2}\X \bbeta + \bm{\Sigma}^{-1/2}\bm{\varepsilon}$. The transformed errors have covariance $\bm{\Sigma}^{-1/2}\bm{\Sigma}\bm{\Sigma}^{-1/2} = \bm{I}$: assumptions (A1)--(A3) hold with design $\tilde{\X} = \bm{\Sigma}^{-1/2}\X$. Applying OLS in the new coordinates and multiplying back gives exactly $\bhat_{\mathrm{GLS}}$; Gauss--Markov makes it BLUE for the transformed problem, and linear unbiasedness is preserved under the invertible transformation.
\end{proof}

\subsection*{Chapter \ref{sec:probability}}

\begin{exercise}
Under (A1)--(A4) with $\X^\top\X = n\bm{I}$ (orthonormal design scaled by $n$), give the exact distribution of $\norm{\bhat - \bbeta}^2$ and use it to build a $1-\alpha$ confidence \emph{ellipsoid} for $\bbeta$.
\end{exercise}
\begin{proof}
By Theorem~\ref{thm:dist-bhat}, $\bhat - \bbeta \sim \mathcal{N}(\zero, \sigma^2 (\X^\top\X)^{-1}) = \mathcal{N}(\zero, \tfrac{\sigma^2}{n}\bm{I})$, so
\[
  \frac{n \norm{\bhat - \bbeta}^2}{\sigma^2} = \sum_{j} \frac{n(\hat\beta_j - \beta_j)^2}{\sigma^2} \sim \chi^2_{p+1}.
\]
Replacing $\sigma^2$ by the independent $\hat\sigma^2 \sim \sigma^2 \chi^2_{n-p-1}/(n-p-1)$ (Theorem~\ref{thm:chi2}) and taking the ratio,
\[
  \frac{n \norm{\bhat - \bbeta}^2}{(p+1)\hat\sigma^2} \sim F_{p+1,\, n-p-1}.
\]
Hence $\big\{ \bbeta : \norm{\bhat - \bbeta}^2 \leq \frac{(p+1)\hat\sigma^2}{n} F_{p+1,\,n-p-1,\,1-\alpha} \big\}$ is an exact $1-\alpha$ confidence ellipsoid (a sphere in these coordinates).
\end{proof}

\subsection*{Chapter \ref{sec:bias-variance}}

\begin{exercise}\label{ex:ridge}
Derive the ridge solution \eqref{eq:ridge-sol} from the penalized objective, and prove that the ridge estimator is biased with $\bias = -\lambda(\X^\top\X + \lambda\bm{I})^{-1}\bbeta$.
\end{exercise}
\begin{proof}
The gradient of $\norm{\y - \X\bbeta}^2 + \lambda\norm{\bbeta}^2$ is $-2\X^\top\y + 2\X^\top\X\bbeta + 2\lambda\bbeta$; setting it to zero and noting $\X^\top\X + \lambda\bm{I} \succ \zero$ gives \eqref{eq:ridge-sol}. Substitute the model:
\[
  \E[\bhat^{\mathrm{ridge}}] = (\X^\top\X + \lambda\bm{I})^{-1}\X^\top\X\,\bbeta
  = \bbeta - \lambda(\X^\top\X + \lambda\bm{I})^{-1}\bbeta,
\]
using $(\X^\top\X + \lambda\bm{I})^{-1}(\X^\top\X) = \bm{I} - \lambda(\X^\top\X+\lambda\bm{I})^{-1}$. The bias term vanishes as $\lambda \to 0$ and saturates as $\lambda \to \infty$ (estimator $\to \zero$).
\end{proof}

\begin{exercise}
In the orthogonal-design ridge setting of Section~7.3, find the $\lambda^\star$ minimizing $\MSE_j$ exactly, and the resulting minimum MSE. For which $\beta_j$ is shrinkage useless?
\end{exercise}
\begin{proof}
$\MSE_j(\lambda) = (\lambda^2\beta_j^2 + \sigma^2)/(1+\lambda)^2$. Setting the derivative to zero:
\[
  \frac{d}{d\lambda}\MSE_j = \frac{2\lambda \beta_j^2 (1+\lambda)^2 - (\lambda^2\beta_j^2 + \sigma^2)\,2(1+\lambda)}{(1+\lambda)^4} = 0
  \;\Longrightarrow\; \lambda^\star = \frac{\sigma^2}{\beta_j^2}.
\]
Then $\MSE_j(\lambda^\star) = \frac{\sigma^2 \beta_j^2}{\beta_j^2 + \sigma^2} < \sigma^2 = \MSE_j(0)$ for any $\beta_j \neq \infty$. As $\beta_j^2 \to \infty$, $\lambda^\star \to 0$ and the gain vanishes: shrinkage is useless (though never harmful in this calculation) for very strong signals — bias--variance thinking allocates shrinkage to weak coefficients.
\end{proof}
